% =====================================================================
%  PASTE-READY PATCHES FOR THE REVIEW RESPONSE
%
%  This file is NOT part of the build. Nothing \input's it. It is a
%  working document: each block below is text to paste into main.tex or
%  supplementary.tex, with the location and the reason recorded beside
%  it, so that a change made in Overleaf can be traced back to the code
%  fact that motivated it.
%
%  Assumes \newcommand{\ours}{SCALAR} is in the preamble (it is, in
%  supplementary.tex line 27). If main.tex uses a different macro,
%  substitute it.
%
%  STATUS KEY
%    [DONE]  already committed to the repo on this branch
%    [PASTE] you must paste this into Overleaf
% =====================================================================


% ---------------------------------------------------------------------
% PATCH 1  [PASTE]   main.tex, Sec. 4.3 "Robustness to Test-Time Missing
%                    Modalities", the \paragraph{Protocol.} block.
%                    Two REPLACEMENTS in place, not an addition.
%
% WHY. The review's most damaging claim was that the headline robustness
% experiment is invalid, on the reasoning that a masked modality reaches
% the Gramian baseline as a zero vector, which puts a zero row AND column
% into G, forces det(G)=0 for every masked clip, and reduces the ranking
% among masked clips to tie-breaking.
%
% The reviewer was not careless. The Protocol paragraph already contains
% BOTH halves of the answer, three lines apart, and they read as a
% contradiction:
%
%   "A masked modality is represented by zeros in the common trunk ..."
%   "... In GRAM, a missing modality contributes a unit factor to the
%    determinant ..."
%
% Zeros cannot produce a unit factor, so a reader who knows Gramians
% believes the first sentence and treats the second as wishful. Both are
% in fact true: the loader writes zeros, and presence is RECOVERED from
% them by a norm test (utils/volume.py:220, present_from_feats, ||z||>0.5)
% before the Gramian is built (utils/volume.py:266-267). That middle step
% appears nowhere in the paper. Adding a sentence elsewhere does not help;
% the two sentences below have to stop contradicting each other.
% ---------------------------------------------------------------------

% 1a. REPLACES: "A masked modality is represented by zeros in the common
%     trunk, consistent with the input produced by a loader for an absent
%     stream."

A masked modality is written as zeros in the common trunk, matching what
a loader produces for an absent stream; presence is then recovered from
the embedding norm before any score is formed, so the zeros are a
transport convention and never reach an aggregation function as a vector.

% 1b. REPLACES: "In GRAM, a missing modality contributes a unit factor to
%     the determinant, resulting in the method's lower-arity formula,
%     which is already implemented at arities 2 to 4 across datasets."
%
%     Keep the sentence that follows it in the paper ("The r=0 setting is
%     byte-identical to the standard protocol ..."): it is correct, and
%     after this replacement it reads as the consequence rather than as an
%     unsupported assurance.

In \textsc{Gram}, the absent axis is made orthonormal rather than
zero-filled: it contributes a factor of exactly one to the determinant,
which therefore reduces to the sub-Gramian over the modalities the clip
retains --- the method's own lower-arity formula, already implemented at
arities 2 to 4 across datasets. Zero-filling instead would place a zero
row \emph{and} column in the Gram matrix and force $\det G = 0$ for
\emph{every} masked clip, collapsing the ranking among them. We do not do
this: the baseline is given a treatment its release does not support,
which discards incomplete clips outright (Sec.~\ref{sec:maskedvol}).


% ---------------------------------------------------------------------
% PATCH 2  [PASTE]   Table 3 caption.
%
% The canonical text lives in scripts/build_missing_table.py
% (MAIN_CAPTION), because everything in tables_final/ must be generator
% output -- scripts/check_tables_generated.py enforces it. Regenerating
% with
%   python3 scripts/build_missing_table.py --all
% puts the text below into
% experiments/results/tables_final/table3_missing.tex.
%
% BUT: the submitted PDF's Table 3 caption reads "For GRAM, the volume is
% computed over the modalities remaining in each clip", which the
% generator has never emitted -- it is a hand-shortened paraphrase. So
% Table 3 in Overleaf is NOT \input from tables_final/, and regenerating
% will not reach it. Either paste the clause below into the Overleaf
% caption, or re-sync the whole table from the generator so the two stop
% drifting. The shortened paraphrase is precisely the version that leaves
% the zero-fill reading open.
% ---------------------------------------------------------------------

% $\star$: authors' released checkpoint. Its release cannot represent a
% missing modality at all and discards such clips, so we extend it
% rather than penalise it: a missing modality is \emph{not} zero-filled
% --- which would place a zero row and column in the Gram matrix and
% force $\det G {=} 0$ for \emph{every} masked clip, collapsing the
% ranking --- but is made an orthonormal phantom axis contributing a
% factor of exactly $1$, so the score reduces to that clip's volume over
% the modalities it retains (Sec.~\ref{sec:maskedvol}).


% ---------------------------------------------------------------------
% PATCH 3  [DONE]   supplementary.tex, new subsection inserted before
%                   \subsection{Two-Stage Results Under Masking}.
%
% Reproduced in full so the whole response reads in one place. Already
% committed; pull rather than paste.
%
% The last paragraph is the load-bearing one. It concedes the correct,
% stronger version of the reviewer's concern -- mixed-arity volumes
% share no common scale, a cosine does -- which costs nothing because it
% explains the gap in Table 3 rather than undermining it.
% ---------------------------------------------------------------------

\subsection{How the Gramian Baseline Sees a Missing Modality}
\label{sec:maskedvol}

The missing-modality comparison is only meaningful if the Gramian
baseline is given a competent way to score an incomplete clip, so we
state the construction explicitly.

The obvious implementation is to zero-fill: substitute $z_m = 0$ for the
absent modality and compute the volume unchanged. This is catastrophic,
and not in an interesting way. A zero vector puts a zero row \emph{and}
a zero column into the Gram matrix $G$, so $\det G = 0$ \emph{exactly},
for every masked clip regardless of its content. Every masked clip
receives the same constant score, the ranking among them is decided by
tie-breaking, and any resulting number would measure the tie-break rule
rather than the aggregator. We do not do this.

Instead, an absent modality is turned into an orthonormal phantom axis.
Writing $p \in \{0,1\}^{|M|+1}$ for the presence indicator (the query is
always present), the masked Gramian is
\begin{equation}
\tilde{G} = G \odot (p\,p^{\!\top}) + \operatorname{diag}(\mathbf{1} - p),
\label{eq:maskedgram}
\end{equation}
which zeroes the absent row and column and then writes a $1$ on its
diagonal. Expanding the determinant along that row gives
$\det \tilde{G} = 1 \cdot \det G_{P}$, where $G_{P}$ is the sub-Gramian
over the present modalities: the absent axis contributes a factor of
exactly one and the clip is scored at its own lower arity, with no
imputed content and no discarded candidate. With $p = \mathbf{1}$ the
construction reproduces the unmasked volume bit-for-bit, so the $r{=}0$
row of the table is the released baseline's own protocol, unmodified.

The same routine and the same presence rule ($\|z_m\| \le 0.5$, since
present embeddings are unit-norm and the loader zero-fills what it
cannot load) are used for training, validation and every evaluation, for
every method. Both identities are asserted in the unit tests: that the
masked volume equals the directly computed lower-arity volume, and that
it is separated from the degenerate zero-fill case.

Two consequences are worth stating plainly. First, this treatment is
\emph{more} favourable than the released implementation supports, which
cannot represent an incomplete clip and discards it; the baseline is
therefore not handicapped by the harness. Second, it leaves a genuine
asymmetry that the scores themselves cannot remove: volumes at different
arities are not on a common scale, so a mixed-arity gallery is ranked by
a quantity whose units vary across candidates, whereas \ours{} returns a
cosine in $[-1,1]$ at every arity. We regard this as part of the
explanation for the gap in Table~3 rather than an artefact of it.


% ---------------------------------------------------------------------
% PATCH 4  [PASTE]   main.tex, Sec. 3, "Cross-arity consistency", the
%                    sentence immediately after Eq. 7.
%
% CORRECTION to an earlier reading of this patch. Eq. 7 is typeset as
%
%     L_mask = 1 - <mu_M, mu_K> + (s_M - sg[s_K])^2
%
% with sg[.] on s_K ONLY, which is exactly what model/losses_sca.py:73-76
% does:
%     if term1: loss += (1.0 - (mu_M * mu_K).sum(-1)).mean()
%     if term2: loss += ((s_M - s_K.detach()) ** 2).mean()
%
% So the EQUATION IS CORRECT AS PRINTED and needs no change. What the
% paper does not do is give any reason for the asymmetry, which is why it
% reads as an inconsistency rather than a decision, and why the review
% treated it as an oversight. This is therefore an addition, not a fix.
%
% REPLACES: "where $\mathrm{sg}[\cdot]$ is the stop-gradient. The first
% term preserves the direction of the representation and the second
% preserves the calibration of the score when a modality is removed."
% ---------------------------------------------------------------------

where $\mathrm{sg}[\cdot]$ is the stop-gradient. The two terms are
gradient-coupled differently, and deliberately. The score term treats the
complete-set score as a fixed reference, so calibration moves the reduced
view onto the complete one rather than letting the two meet at some
intermediate value that neither view would produce at test time. The
direction term carries gradient through both centroids: agreement across
arities is a property the representation must have, not a target that one
view holds for the other, and freezing $\mu_{\mathcal{K}}$ would exempt
the complete-modality view from the constraint it imposes.


% ---------------------------------------------------------------------
% PATCH 5  [PASTE]   main.tex, Sec. 3, "Graded semantic supervision".
%                    TWO edits, one before Eq. 8 and one after it.
%
% WHY. sharpen_targets (model/losses_sca.py:28) is
%     F.softmax(s_star / tau_star, dim=-1)
% over ALL entries, so an entry sparsified away receives logit 0 and is
% treated as a value. cal_known_only (losses_sca.py:55-57) masks the
% CALIBRATION term only, not the KL term. The paper's "an absent entry
% means unknown rather than zero" is stated generally but is true of half
% the loss.
%
% Both edits are needed. 5a alone leaves the overclaim standing in weaker
% form; 5b alone contradicts the sentence above it.
% ---------------------------------------------------------------------

% 5a. REPLACES: "Let $\Omega = \{(i,j) : S^\star_{ij} > 0\}$ denote the
%     pairs with an observed target; the cache is sparsified, so an
%     absent entry means unknown rather than zero."

Let $\Omega = \{(i,j) : S^{\star}_{ij} > 0\}$ denote the pairs with an
observed target. The cache is sparsified, so an absent entry means the
affinity is unknown rather than zero, and the calibration term below is
therefore restricted to $\Omega$.

% 5b. APPENDS after: "where the rescaling $2S^\star - 1$ maps affinities
%     onto the cosine range $[-1,1]$."

The KL term does not admit the same restriction: it is normalised over
the full row, and renormalising over a support that varies per row would
make the objective incomparable across rows. Absent entries there
consequently enter at the softmax's uniform baseline weight. The ranking
term thus treats an unknown affinity as unremarkable, while the
calibration term abstains on it.


% ---------------------------------------------------------------------
% PATCH 6  [DONE]   Anonymity. Two leaks, both fixed on the branch.
%
%  (a) supplementary.tex, the evaluation-environment audit table, had a
%      row naming the third-party reimplementation by project name. It
%      now reads:
%
%          GRAM, released checkpoint & third-party reimplementation & 53.4 \\
%
%      If you have edited that table in Overleaf since, re-check the
%      second column of the second row.
%
%  (b) experiments/results/tables_final/trainable_params.json recorded
%      absolute cluster paths, which embed the allocation account. The
%      committed entries are rewritten and scripts/count_trainable.py
%      now anonymises new ones via anon_source(). Not in the PDF, but it
%      is in the repo, and a repo link in the submission would carry it.
% ---------------------------------------------------------------------


% ---------------------------------------------------------------------
% STILL OPEN -- both need cluster time, neither is a paste
%
%  (i) MAX-OVER-MODALITIES BASELINE ROW. The review asks for it and the
%      ask is fair: the query-weighted centroid does tend to argmax as
%      $\tau_w \to 0$, and tests/test_query_centroid.py proves the limit
%      exactly. It is cheap -- $\tau = 10^{-5}$ through
%      scripts/try_query_centroid.py on the cached features, no
%      retraining. Our own committed sweep already refutes the strong
%      form of the claim: at $\tau{=}0.01$ (effectively the max) R@1 is
%      at or below $\tau{=}0.1$ on 9 of 9 feature dumps, by up to 9.6
%      points, and on AudioCaps the max limit loses 6--12 points against
%      the best single modality. Both endpoints of the temperature range
%      are worse than the operating point, which is the strongest form
%      of the answer -- but it should be a table row, not a sweep buried
%      in the harvest.
%
% (ii) 54.6 vs 54.8 LABELLING. Three-seed mean against the seed-0
%      reference, never labelled as such. One clause in the table note.
% ---------------------------------------------------------------------
